% TOP_PROBLEM: {"id":30007725,"problem_number":"LOCAL-30007725","title":"Provider payment and patient welfare","category_id":21,"category":{"id":21,"name":"economics","display_name":"Economics","description":"Open economic theory statements, published research agendas, and proposed scoped specifications; question provenance and historical priority are qualified per record.","slug":"economics","order_index":21,"created_at":"2026-10-08T00:00:00Z"},"status":"research_question","external_url":null,"legacy_ids":[],"published":false,"statement":"Estimate whether a specified primary-care performance-payment contract improves three-year patient welfare, including unrewarded outcomes and patient burden.","economics":{"original_id":214,"field":"Health economics","kind":"empirical","formalization_basis":"adapted_research_question","source_status":"research_agenda","priority_status":"earliest_found","priority_note":"Earliest explicit statement located in this audit; no exhaustive priority search or first-ever claim. The mathematical specification below is adapted, not quoted from the source.","earliest_verified_reference":{"authors":["Martin Gaynor","Kate Ho","Robert J. Town"],"title":"The Industrial Organization of Health Care Markets","year":2014,"url":"https://kateho.scholar.princeton.edu/sites/g/files/toruqf4136/files/kateho/files/healthcare_io_3_24_14.pdf","doi":"10.3386/w19800","locator":"§7 Models of Health Plan Choice, p. 40; §8, p. 43; §9 Future Directions, pp. 47–48","evidence_summary":"The authors explicitly discuss unanswered questions about consumer inertia, adverse selection, provider incentives, and consumer search frictions."},"current_reference":{"authors":["Anthony Scott","Peter Sivey","Driss Ait Ouakrim","Laura Willenberg","Lena Naccarella","John Furler","Doris Young"],"title":"The effect of financial incentives on the quality of health care provided by primary care physicians","year":2022,"url":"https://www.cochrane.org/evidence/CD008451_effect-financial-incentives-quality-health-care-provided-primary-care-physicians","doi":"10.1002/14651858.CD008451.pub2","locator":"Authors’ conclusions, displayed publication date 1 May 2022","evidence_summary":"The review requests evidence on broader outcomes, unintended consequences, subgroups, and comparative costs and effects."},"formalization_note":"The source explicitly identifies the underlying gap or agenda. Population, horizon, interventions, estimands, and auxiliary assumptions are proposed operational choices, not a canonical source theorem. Partial later answers are compatible with this scoped question; the citation alone does not prove absence of every subsequent solution."},"statusQualification":"Published question or research agenda verified at the cited locator. The mathematical specification is adapted for this catalogue and includes additional assumptions; source publication does not establish first-ever priority or certify this exact model remains unsolved. The source explicitly identifies the underlying gap or agenda. Population, horizon, interventions, estimands, and auxiliary assumptions are proposed operational choices, not a canonical source theorem. Partial later answers are compatible with this scoped question; the citation alone does not prove absence of every subsequent solution.","statusReviewedAt":"2026-10-08"}
% FIRST_POSTED: 2026-10-08
% ENTRY_KIND: statement_only
% !TeX program = lualatex
\documentclass[11pt]{article}
\usepackage{fontspec}
\usepackage[a4paper,margin=25mm]{geometry}
\usepackage{amsmath,amssymb,mathtools}
\usepackage{xurl}
\usepackage[hidelinks]{hyperref}
\newcommand{\sref}[2]{\hyperlink{source:#1:#2}{[S#2]}}
\setlength{\emergencystretch}{3em}
\begin{document}
% problemId:30007725
\section*{Provider payment and patient welfare}\label{30007725}
\subsection{Definitions and mathematical statement}
Consider primary-care practices serving adults with diabetes in one national health system. Randomize practices to \(a=1\), a specified performance-payment contract, or \(a=0\), their existing payment arrangement, for three years. Let \(Q_i(a)\) denote patient quality-adjusted life-years, \(C_i(a)\) total treatment resource costs, and \(B_i(a)\) appointment and administrative burden. Define
\[\tau=E[\lambda\{Q_i(1)-Q_i(0)\}-\{C_i(1)-C_i(0)\}-\rho\{B_i(1)-B_i(0)\}],\]
where \(\lambda>0\) and \(\rho>0\) are prespecified monetary valuations, and the expectation covers all baseline eligible patients. Measure untreated conditions and avoidable hospitalizations alongside the contract’s rewarded indicators. Under consistency and randomized practice assignment, estimate intention-to-treat effects, adjusting inference for practice clustering. Address spillovers between practices or define a corresponding exposure model. Determine whether improved measured performance yields net patient benefit, and whether effects differ by baseline disadvantage. Transfers to providers enter fiscal accounting, whereas welfare uses their actual resource opportunity cost. This adapts explicit source concerns about incentive effects on quality and unintended outcomes; it does not presume that every performance-pay contract has the same effect.

\par\textbf{Formulation and scope.} The source explicitly identifies the underlying gap or agenda. Population, horizon, interventions, estimands, and auxiliary assumptions are proposed operational choices, not a canonical source theorem. Partial later answers are compatible with this scoped question; the citation alone does not prove absence of every subsequent solution.

\par\textbf{Open-question provenance.} An explicit question or research agenda was verified at the source locator. This does not by itself certify that every later version remains unresolved \sref{30007725}{1}.

\par\textbf{Historical priority.} Earliest explicit statement located in this audit; no exhaustive priority search or first-ever claim. The mathematical specification below is adapted, not quoted from the source.

\subsection{Short English statement}
Estimate whether a specified primary-care performance-payment contract improves three-year patient welfare, including unrewarded outcomes and patient burden.

\subsection{Sources}
\begin{itemize}\raggedright
\item[\textnormal{[S1]}]\hypertarget{source:30007725:1}{} Martin Gaynor, Kate Ho, Robert J. Town. 2014. The Industrial Organization of Health Care Markets. §7 Models of Health Plan Choice, p. 40; §8, p. 43; §9 Future Directions, pp. 47–48.
\url{https://kateho.scholar.princeton.edu/sites/g/files/toruqf4136/files/kateho/files/healthcare_io_3_24_14.pdf}.
DOI: 10.3386/w19800.
{\small\textit{Earliest verified question/agenda reference; historical priority qualified below. The authors explicitly discuss unanswered questions about consumer inertia, adverse selection, provider incentives, and consumer search frictions.}}

\item[\textnormal{[S2]}]\hypertarget{source:30007725:2}{} Anthony Scott, Peter Sivey, Driss Ait Ouakrim, Laura Willenberg, Lena Naccarella, John Furler, Doris Young. 2022. The effect of financial incentives on the quality of health care provided by primary care physicians. Authors’ conclusions, displayed publication date 1 May 2022.
\url{https://www.cochrane.org/evidence/CD008451_effect-financial-incentives-quality-health-care-provided-primary-care-physicians}.
DOI: 10.1002/14651858.CD008451.pub2.
{\small\textit{Supporting or current literature; see evidence qualification. The review requests evidence on broader outcomes, unintended consequences, subgroups, and comparative costs and effects.}}
\end{itemize}
\end{document}
